\chapter{Le Processus de Gestion d'Incident (ISO 27035)}
\label{chap:procincident}

Pour gérer l'imprévisible, il faut une méthodologie rigoureuse. La norme \textbf{ISO 27035} fournit le cadre international pour gérer les incidents de sécurité de l'information de manière structurée.

\section{Présentation de l'ISO 27035}

\subsection{Objectif de la Norme}

\begin{boxdef}
\textbf{ISO 27035} : Norme internationale intitulée "Gestion des incidents de sécurité de l'information". Elle fournit des lignes directrices pour planifier, préparer, détecter, signaler et répondre aux incidents.
\end{boxdef}

L'objectif n'est pas seulement de réagir, mais d'apprendre de chaque incident pour renforcer le Système de Management de la Sécurité (SMSI). Elle s'inscrit dans le cycle PDCA (Act).

\subsection{Les 5 Phases du Processus}

La norme structure la gestion d'incident en cinq phases séquentielles :

\begin{enumerate}
    \item \textbf{Planifier :} Définir la politique d'incident.
    \item \textbf{Préparer :} Mettre en place les équipes et les outils.
    \item \textbf{Détecter et Signaler :} Identifier les événements suspects.
    \item \textbf{Évaluer et Réagir :} Traiter l'incident.
    \item \textbf{Apprendre :} Améliorer le système.
\end{enumerate}

% ------------------------------------------------------------------------------
\section{Phase 1 et 2 : Planification et Préparation}
% ------------------------------------------------------------------------------

\subsection{Politique de Gestion d'Incident}

L'entreprise doit définir clairement :
\begin{itemize}
    \item Ce qui constitue un incident.
    \item Qui est responsable de la réaction.
    \item Les délais de remontée d'information.
\end{itemize}

\subsection{Les Équipes d'Intervention (CSIRT / CERT)}

La préparation implique la création d'une équipe dédiée.

\begin{boxdef}
\begin{itemize}
    \item \textbf{CERT (Computer Emergency Response Team) :} Terme historique, souvent utilisé pour les équipes nationales ou sectorielles.
    \item \textbf{CSIRT (Computer Security Incident Response Team) :} Terme standard pour une équipe interne à l'entreprise.
\end{itemize}
\end{boxdef}

Le CSIRT est le point de contact unique pour la gestion des incidents. Il a pour mission : Réception, Tri, Analyse, Réponse, Coordination.

\begin{boxlizbiz}
\textbf{Organisation chez LiZbiZ :}
LiZbiZ n'a pas encore de CSIRT dédié. La DSI a demandé à un groupe de 3 techniciens de former une "Cellule Cyber" de garde. L'un d'eux est désigné comme "Leader Incident".
\end{boxlizbiz}

% ------------------------------------------------------------------------------
\section{Phase 3 : Détection et Signaler}
% ------------------------------------------------------------------------------

C'est le moment où l'on passe de l'événement brut à l'incident confirmé.

\subsection{Les Outils de Détection (SOC et SIEM)}

\begin{boxdef}
\textbf{SOC (Security Operations Center) :} Structure organisationnelle (une équipe) qui surveille le SI 24/7 pour détecter les anomalies.
\end{boxdef}

\begin{boxdef}
\textbf{SIEM (Security Information and Event Management) :} Outil logiciel qui centralise les logs, corrèle les événements et génère des alertes.
\end{boxdef}

Le SIEM est l'outil technique indispensable. Il collecte les logs des pare-feux, des serveurs, des antivirus et cherche des motifs d'attaque.

\subsection{Tri et Qualification}

Toutes les alertes ne sont pas des incidents. Le processus de tri (Triage) doit répondre à :
\begin{itemize}
    \item L'alerte est-elle un Faux Positif ?
    \item Quelle est la gravité réelle ?
    \item Quelle est l'urgence de traitement ?
\end{itemize}

% ------------------------------------------------------------------------------
\section{Phase 4 : Réaction et Traitement}
% ------------------------------------------------------------------------------

Une fois l'incident confirmé, la réaction suit un cycle strict.

\subsection{Confinement (Containment)}

C'est la priorité absolue. Limiter les dégâts avant de nettoyer.
\begin{itemize}
    \item Isoler la machine du réseau (VLAN de quarantaine).
    \item Bloquer une IP externe au pare-feu.
    \item Désactiver un compte utilisateur compromis.
\end{itemize}

\subsection{Éradication et Rétablissement}

\begin{itemize}
    \item \textbf{Éradication :} Supprimer le malware, reboucher la faille.
    \item \textbf{Rétablissement :} Restaurer les données, remettre le service en ligne.
\end{itemize}

\subsection{Les Outils Procéduraux : Playbook et Runbook}

Pour réagir vite, on n'improvise pas. On suit des scripts.

\begin{boxdef}
\begin{itemize}
    \item \textbf{Playbook :} Document qui décrit la procédure logique de réponse à un type d'incident donné (ex: "Procédure en cas de Ransomware"). C'est le "Quoi faire".
    \item \textbf{Runbook :} Document technique détaillant les commandes exactes et les clics à effectuer (ex: "Commande Linux pour isoler une interface"). C'est le "Comment faire".
\end{itemize}
\end{boxdef}

% ------------------------------------------------------------------------------
\section{Cas LiZbiZ : Simulation de Détection dans le Lab}
% ------------------------------------------------------------------------------

\subsection{Scénario : L'Alerte Nocturne}

\begin{boxlizbiz}
\textbf{Contexte :} Il est 3h du matin. Le serveur SIEM du Lab GNS3 génère une alerte critique.
\end{boxlizbiz}

\subsection{Analyse de l'Alerte (Travaux Pratiques)}

Les étudiants se connectent à l'interface du SIEM (ex: Splunk ou ELK Stack simulé).

\textbf{L'alerte indique :}
\begin{itemize}
    \item \textbf{Source :} VLAN 40 (IoT Logistique).
    \item \textbf{Destination :} VLAN 20 (Serveur ERP).
    \item \textbf{Type :} Tentative de connexion SSH multiple (Brute Force).
    \item \textbf{État :} Succès (Authentification réussie).
    \item \textbf{Utilisateur :} admin\_erp.
\end{itemize}

\textbf{Questions pour l'analyste (Étudiant) :}
\begin{enumerate}
    \item \textbf{Qualification :} Est-ce un incident ou un événement normal ? (Justifiez : un scanner IoT ne doit pas se connecter en SSH à l'ERP).
    \item \textbf{Confinement :} Quelle action technique immédiate proposez-vous sur le Firewall pFsense pour stopper l'attaque ? (Proposition : Bloquer l'IP source ou couper le flux VLAN 40 -> 20).
    \item \textbf{Consultation du Playbook :} Le Playbook "Accès non autorisé" demande de vérifier si d'autres machines sont touchées. Quelle commande Linux utiliseriez-vous sur le serveur ERP pour voir qui est connecté ? (Proposition : \texttt{who} ou \texttt{last}).
\end{enumerate}

% ==============================================================================
% Fichier : 04_module_incident.tex (Extrait Chapitre 3)
% Description : Investigation Numérique (Forensics)
% ==============================================================================